%! Unit Opener: Why Algebra Fluency Matters — Unit 1, Lesson 1.0 — Group Activity
\documentclass[10pt]{article}
\usepackage{atda-article}
\usepackage{atda-boxes}
\newcommand{\TallMath}[1]{$\displaystyle #1\rule[-1.4em]{0pt}{3.2em}$}
\setlength{\workrowsep}{6pt}

\begin{document}

\pageheader{Unit 1, Lesson 1.0}{Group Activity --- The Spirit-Wear Sale}
% NAMESTRIP: no \namedateperiod here. The student writes their name once, on the
% cover this component is stapled behind. See references/conventions.md.

\vspace{0.15in}

\begin{scenariobox}[The Scenario]{royal}
\small
The robotics club is pricing T-shirts for the fall sale. Their survey says that if they charge
$x$ dollars per shirt, they will sell about $120-4x$ shirts. So the money the club takes in ---
its \textbf{revenue}, in dollars --- is
\[ R(x) = x(120-4x). \]
Work with your group. \textbf{Everyone starts at Tier R.} When your whole group can do Tier R
without help, move on.
\end{scenariobox}

\vspace{0.10in}

\boxguard[20]
\begin{tcolorbox}[colback=white, colframe=black!40, title=\textbf{Tier R --- Remediate},
    fonttitle=\bfseries, arc=2mm, left=3mm, right=3mm, top=2mm, bottom=2mm, breakable]
\small
Warm your hands up on the two moves the rest of the activity needs.

\begin{enumerate}[leftmargin=1.6em, itemsep=4pt, topsep=4pt]
  \item Multiply and write in standard form: $(x+3)(x+5)$
\begin{work}
  (x+3)(x+5) &= x^2+5x+3x+15 \\
             &= x^2+8x+15
\end{work}

  \item Distribute: $2x(x-6)$
\begin{work}
  2x(x-6) &= 2x^2-12x
\end{work}

  \item Factor out the greatest common factor: $15y^3-25y^2$
\begin{work}
  15y^3-25y^2 &= 5y^2(3y)-5y^2(5) \\
              &= 5y^2(3y-5)
\end{work}
\end{enumerate}
\end{tcolorbox}

\vspace{0.10in}

\boxguard[20]
\begin{tcolorbox}[colback=white, colframe=black!40,
    title=\textbf{Tier A --- Approaching Proficiency},
    fonttitle=\bfseries, arc=2mm, left=3mm, right=3mm, top=2mm, bottom=2mm, breakable]
\small
Now the club's revenue model, $R(x)=x(120-4x)$.

\begin{enumerate}[leftmargin=1.6em, itemsep=4pt, topsep=4pt]
  \item Expand $R(x)$ and write it in standard form.
\begin{work}
  R(x) &= x(120-4x) \\
       &= 120x-4x^2 \\
       &= -4x^2+120x
\end{work}

  \item The factored form has two zeros. Find them.
\begin{work}
  x &= 0 \text{ or } 120-4x = 0 \\
  x &= 0 \text{ or } x = 30
\end{work}

  \item Each zero is a price at which the club takes in \emph{no money}. Explain, in context, why
        each one brings in nothing --- the two reasons are different.

\writelines{2}

\tcbbreak
  \item Find $R(10)$ and write one sentence saying what it means for the club.
\begin{work}
  R(10) &= 10\big(120-4(10)\big) \\
        &= 10(80) \\
        &= 800
\end{work}
\writeline

  \item Your treasurer asks, ``At what prices do we make nothing at all?'' Which form of $R$
        answers that fastest, and why?

\writelines{2}
\end{enumerate}
\end{tcolorbox}

\vspace{0.10in}

\boxguard[22]
\begin{tcolorbox}[colback=white, colframe=black!40, title=\textbf{Tier E --- Extension},
    fonttitle=\bfseries, arc=2mm, left=3mm, right=3mm, top=2mm, bottom=2mm, breakable]
\small

\begin{enumerate}[leftmargin=1.6em, itemsep=4pt, topsep=4pt]
  \item \textbf{Critique the work.} A student expanded a different club's model like this:
        \[ -4(x-5)(x+2) \;=\; -4\big(x^2-3x-10\big) \;=\; -4x^2-3x-10. \]
        Exactly one step is wrong. Name the error, then write the correct expansion.
\begin{work}
  -4(x-5)(x+2) &= -4\big(x^2-3x-10\big) \\
               &= -4x^2+12x+40
\end{work}
\writeline

  \item \textbf{Use the symmetry.} A parabola is symmetric, so the price that earns the
        \emph{most} revenue sits exactly halfway between the two zeros of $R(x)=x(120-4x)$.
        Find that price and the revenue it brings in.
\begin{work}
  x &= \dfrac{0+30}{2} = 15 \\
  R(15) &= 15\big(120-4(15)\big) \\
        &= 15(60) = 900
\end{work}

  \item \textbf{Justify.} A club member argues, ``If we want more money, we should charge more.''
        Use your answer to \#2 to explain why that reasoning fails here.

\writelines{2}
\end{enumerate}
\end{tcolorbox}

\end{document}
